\section{Conclusion}
\label{sec:conclusion}

AI agents have successfully carried out a complete thrust measurement archived ALEPH 1994 data at $\sqrt{s}=91.2$~GeV under physicist guidance and collaboration. The analysis propagates uncertainty components arising from statistical, experimental, and theoretical sources and combines them in quadrature. The corresponding correlation matrices are derived from the covariance structure. The implementation of the ISR correction, following the prescription of both published ALEPH event-shape analyses~\cite{Barate:1996fi,Heister:2003aj}, brings the measured thrust distribution into good agreement with the published ALEPH 2004 result, with the only notable
discrepancy at the $T\to 1$ dijet endpoint. The AI-agent operating model demonstrated in this cycle, executable-workflow orchestration under physicist supervision with full
provenance tracking, provides a scalable framework for precision QCD as a testing ground for closed theory--experiment loops in fundamental physics.

\subsection*{Lessons from AI--Physicist Collaboration}

Carrying out this analysis has also produced practical insights into how AI agents and physicists can work together effectively. Our workflow followed a constrained human-in-the-loop cycle in which the physicist defined physics objectives, acceptance criteria, and interpretation authority, while the AI agent owned implementation and execution. Crucially, experienced physicists contribute domain context that cannot yet be reduced to formal specifications alone, including judgments about physically meaningful binning, acceptable regularization behavior, uncertainty communication, and when apparent anomalies warranted deeper investigation. A major future goal identified in this work is to translate this expert knowledge into explicit, code-based verification layers. 

From this experience, several recurring challenge classes emerged. Implicit physics priors, expert expectations about acceptable binning, regularization behavior, or uncertainty presentation, were often assumed rather than explicitly stated, requiring iterative clarification. Figure formatting and style choices (label offsets, legend placement, pad geometry) required many short feedback loops, with cosmetic shortcomings sometimes affecting interpretability and therefore physics conclusions. Synchronization across rapidly evolving scripts, outputs, and note text demanded continuous provenance refreshing. These lessons motivate creation of reusable tools, stronger configuration contracts, clearer channels for encoding expert heuristics into machine readable constraints in future iterations of this program, and new verification layers throughout the autonomous analysis flow.